\chapter[Governing Equations]{Governing Equations in Mantle Dynamics}
\label{equations}
\section{Mechanics of Mantle Dynamics}
Fluid dynamics of the mixture of two phase flow has been extensively discussed.
I will disscuss the three governing equations of mass and momentum referring
to the similar discussion in \cite{AveragedEquationsTwoPhase}, 
\cite{MagmaPhysics} and \cite{DynamicsPartiallyMoltenRock}.
Considering a conservation equation for given volumetric variable $\gamma$ in 
a representative volume element $V$.
\begin{equation}\label{eq:generalconservationlaw}
		  \frac{d}{dt} \int_{V} \gamma d\bold{v} = 
		  - \int_{\partial V} \bold{f} \cdot d\sigma + \int_{V} S d\bold{v}
\end{equation}
where $\bold{f}$ represents flux travelling across the boundary of a representative 
volume element $V$. The minus sign indicates that the positive directions aligns with
normal unit vector pointing outwards. $S$, source/sink term, represents the volumetric
production rate of $\gamma$.
We can apply Gauss's theorem to obtain differential form,
\begin{equation}\label{eq:conservationlawpde}
\frac{\partial \gamma}{\partial t} + \nabla \cdot \bold{f} = S
\end{equation}
We also define the quantity of phase fraction rigorously with the help of an 
indicator function admitting that multiple phase cannot coexist at the same
spatial point, which are regarded as immiscible phases.
We define indicator function,
\begin{equation}\label{eq:indicatorfunction}
		  i(\bold{x}) = \left \{ 
		  \begin{aligned}
					 &1 \quad \text{phase}_i\\
					 &0 \quad \text{not phase}_i
		  \end{aligned}
		  \right.
\end{equation}
We then define macroscopic phase fraction function:
\begin{equation}\label{eq:porosityfunction}
		  \phi_i = \frac{1}{\delta \bold{v}} 
		  \int_{\delta \bold{v}} i(\bold{x}) d\bold{v}
\end{equation}
In the beginning of my discussion, I will discuss two immiscible phases liquid magma 
and solid mantle. Phase fraction of liquid magam is labelled as $\phi_f$ and phase
fraction of solid mantle is labelled as $\phi_s$, where $\phi_f + \phi_s = 1$ is
always attained.

Substituting bulk density $\overline{\rho} = \phi_f \rho_f + \phi_s \rho_s$ for generalized 
volumetric variable $\gamma$ in equation \ref{eq:generalconservationlaw}, I can obtain
conservation law for bulk mass,
\begin{equation}\label{eq:massconservation}
\frac{\partial \overline{\rho}}{\partial t} + \nabla \cdot \overline{\rho \bold{v}} = 0
\end{equation}
with the help of Boussinesq approximation ignoring compressibility of each individual
phase, and extend Boussinesq approximation further applying to density difference 
between each phase except body force term. I can write continuous function for two 
phase mixture simply as,
\begin{equation}\label{eq:continuousfunction}
\nabla \cdot \overline{\bold{v}} = 0
\end{equation}
I then write conservation of momentum with phase averaged term,
\begin{equation}\label{eq:momentumconservation}
\frac{\partial \overline{\rho \bold{v}}}{\partial t} + \nabla \cdot \overline{\rho \bold{v} \otimes \bold{v}}  = 
\overline{\rho}\bold{g} + \nabla \cdot \overline{\bold{\tau}} - \nabla \overline{P}
\end{equation}
where $\bold{\tau}$ represents deviatoric stress,
Ignoring change of inertia term due to near zero Reynolds number in magma and mantle, 
\begin{equation}\label{eq:staticmomentumconservation}
\overline{\rho}\bold{g} + \nabla \cdot \overline{\bold{\tau}} - \nabla \overline{P} = 0
\end{equation}

\vspace{1cm}
Insert more details later!
\vspace{1cm}

\subsection{A Darcy-Stokes Flow System}
Compbining equations ...., I finally arive with the full flow mechanical system of a two-phase system. 


\section{Conservation of species and energy}

\section{Eutectic Phase Behavior}

\section{Rescaling and Nondimensionalization}
